\section*{Wstęp}
\addcontentsline{toc}{section}{Wstęp}

Współczesna inżynierija oprogramowania kładzie szczególny nacisk na modułowość kodu, poprawność alokacji dynamicznej oraz rzetelne prowadzenie dokumentacji technicznej i wersji projektowych. Celem niniejszego sprawozdania jest przedstawienie procesu projektowania, implementacji oraz dokumentowania zaawansowanej struktury danych -- dwukierunkowej listy wiązanej działającej na stercie w języku C++.

Projekt został zrealizowany w ramach zajęć laboratoryjnych na kierunku inżynieryjnym w Akademii Nauk Stosowanych w Nowym Sączu. Obejmuje on nie tylko poprawną implementację generycznych klas węzła, listy oraz iteratora, ale również wykorzystanie wzorców projektowych (\textit{Factory} i \textit{Iterator}), automatyczne generowanie dokumentacji za pomocą narzędzia \textbf{Doxygen} oraz ścisłe przestrzeganie procedur wersjonowania kodu źródłowego przy użyciu systemu \textbf{Git} i platformy GitHub.